\begin{figure*}[htbp]
    \centering
    \includegraphics[width=\linewidth]{../experiments_v2/02_nonstationary_k3_drift/results/gamma_trajectory.pdf}
    \caption{Adaptive forgetting factor $\gamma_t$ trajectory under
    reward shift (left) and cost shift (right)
    ($K{=}3$; 40~seeds, mean $\pm$1\,SD shading).
    %
    The adaptive mechanism uses a dual-EMA residual
    estimator~\cite{cavenaghi2021nonstationary,bifet2007adwin}
    that monitors each arm's absolute prediction residual
    $|r_t - \hat\theta_a^\top x_t|$ via fast
    ($\alpha_s{=}0.3$) and slow ($\alpha_l{=}0.05$) exponential
    moving averages.
    When the fast EMA exceeds the slow EMA beyond a self-calibrated
    noise margin (Eq.~\ref{eq:shift_estimate}), $\gamma_t$ is
    driven from $1.0$ toward the floor $0.999$
    (Eq.~\ref{eq:adaptive_gamma}); otherwise it remains at~$1.0$.
    Horizontal reference lines show the fixed-$\gamma$ baselines.
    %
    \textbf{Left (reward shift):}
    $\gamma_t$ drops sharply near the phase boundary
    (step~${\sim}950$) as prediction residuals spike for
    the swapped arms, then gradually recovers as the reward
    model adapts.  During Phase~1, a minority of seeds
    (${\sim}15$--$25\%$) transiently trigger $\gamma{=}0.999$
    due to burn-in calibration noise, but the population mean
    stays above $0.9995$---confirming the noise margin suppresses
    most stationary-regime fluctuations.
    %
    \textbf{Right (cost shift):}
    $\gamma_t$ remains near $1.0$ throughout the entire
    trajectory, including after the Gemini price drop at
    step~893.  Because reward distributions did not change,
    prediction residuals stay within the noise margin and the
    adaptive mechanism correctly avoids triggering forgetting.
    The small residual variance ($\pm 0.0003$) reflects the
    minority of seeds where arm-switching noise briefly
    exceeds the margin.
    %
    This asymmetric response---active forgetting under reward shift,
    quiescent under cost shift---is the key property that enables
    the adaptive mechanism to match fixed $\gamma{=}0.999$ when
    forgetting is needed ($p{=}0.98$, n.s.) and significantly
    outperform it when forgetting is harmful
    ($\Delta{=}{-}11.1$, $p < 10^{-12}$;
    Table in Section~\ref{sec:nonstat_significance}).
    }
    \label{fig:gamma_trajectory}
\end{figure*}
